A set of runs from one store during physics quality RHIC Run 2024 \auau collision period were used for this analysis. 
These runs were selected based on high calorimeter and global detector data quality.
The calorimeters had a limited set of bad channels for runs in this store (>98\% in the EMCal and >99\% in the HCals).
The global detectors (MBD + ZDC) also had high data quality and were used for centrality determination for each run in the store. 
Data was collected using a hardware level minimum bias trigger requiring at least two PMTs on each side of the MBD detector to fire.
A set of minimum bias event offline selection criteria was also applied to the data to remove events consistent with beam-related backgrounds and non-hadronic interactions.
More details on the minimum bias trigger and event selection used for this analysis is outlined in Section~\ref{sec:detdeta_selection}.

Three MC samples are used for comparison to data and for deriving correction factors mapping the reconstructed \detdeta distributions back to the truth particle level \detdeta. 
1 million minimum bias \auau events at \snn = 200 GeV were used from each generator, \verb|HIJING|~\cite{hijing}, \verb|AMPT|~\cite{ampt}, and \verb|EPOS4|~\cite{epos}.
These events were then passed through the full sPHENIX detector \textsc{GEANT-4}~\cite{GEANT4} simulation and reconstruction to produce simulation datasets that can be compared directly to data.
The MC samples are also reweighted to match the z-vertex distribution in data as well as the particle composition in previous \auau measurements at \snn = 200 GeV. 
More details on the MC samples reweighting for this analysis is included in Section~\ref{sec:detdeta_mcsamples}.